% Part: proof-theory
% Chapter: normalization
% Section: introduction

\documentclass[../../../include/open-logic-section]{subfiles}

\begin{document}

\olfileid[id]{pt}{nor}{seg}

\olsection{Segmen dan Potongan}

!!^a{introduction} yang diikuti !!a{elimination} merupakan jalan memutar
dalam !!a{proof}, dan menormalisasi !!a{proof} berarti menghapus jalan
memutar seperti itu. Namun, aturan \Elim{\lor} dan \Elim{\lexists} bahkan
membuat definisi ``jalan memutar yang ingin kita hapus'' sedikit rumit.
Masalahnya sebagai berikut. Karena aturan-aturan tersebut, dalam suatu jalan
memutar tidak ada jaminan bahwa !!{elimination} \emph{langsung} mengikuti
!!{introduction}. Di antara keduanya dapat terdapat \Elim{\lor} atau
\Elim{\lexists}, misalnya,
\[
\AxiomC{}
\DeduceC{$\lexists[x][!C(x)]$}
\AxiomC{$\Discharge{!A}{x}$}
\DeduceC{$!B$}
\DischargeRule{\Intro{\lif}}{x}
\UnaryInfC{$!A \lif !B$}
\DischargeRule{\Elim{\lexists}}{y}
\BinaryInfC{$!A \lif !B$}
\AxiomC{}
\DeduceC{$!A$}
\RightLabel{\Elim{\lif}}
\BinaryInfC{$!B$}
\DisplayProof
\]
Bahkan, sejumlah berapa pun \Elim{\lor} atau \Elim{\lexists} dapat
memisahkan \Intro{\lif} dari \Elim{\lif}, misalnya,
\[
\AxiomC{}
\DeduceC{$\lexists[x][!C(x)]$}
\AxiomC{}
\DeduceC{$!D \lor !E$}
\AxiomC{$\Discharge{!A}{x}$}
\DeduceC{$!B$}
\DischargeRule{\Intro{\lif}}{x}
\UnaryInfC{$!A \lif !B$}
\AxiomC{$\Discharge{!A}{x}$}
\DeduceC{$!B$}
\DischargeRule{\Intro{\lif}}{x}
\UnaryInfC{$!A \lif !B$}
\RightLabel{\Elim{\lor}}
\TrinaryInfC{$!A \lif !B$}
\DischargeRule{\Elim{\lexists}}{y}
\BinaryInfC{$!A \lif !B$}
\AxiomC{}
\DeduceC{$!A$}
\RightLabel{\Elim{\lif}}
\BinaryInfC{$!B$}
\DisplayProof
\]
Contoh ini menunjukkan bahwa \Elim{\lif} yang sama bahkan dapat terlibat
dalam beberapa jalan memutar karena mengikuti lebih dari satu~\Intro{\lif}.
Untuk mencakup semua kemungkinan tersebut, kita mendefinisikan jalan
memutar dengan menggunakan jenis tertentu dari barisan kemunculan
!!{formula} dalam bukti~\Log{N1i}, yang disebut \emph{segmen}. Dalam contoh
kita, barisan itu terdiri atas kemunculan~$!A \lif !B$, dengan premis minor
$\Elim{\lor}$ atau $\Elim{\lexists}$ diikuti oleh kesimpulannya. Dalam
contoh tersebut, terdapat dua segmen seperti itu, dan masing-masing memuat
tiga kemunculan~$!A \lif !B$. Berikut definisi resminya.

\begin{defn}
Suatu \emph{segmen} dalam !!{proof}~\Log{N1i}~$\delta$ adalah barisan
kemunculan $!A_1$, \dots, $!A_n$ dari !!{formula}~$!A$ yang sama dalam
$\delta$, dengan ketentuan
\begin{enumerate}
  \item $!A_1$ bukan kesimpulan \Elim{\lor} atau \Elim{\lexists};
  \item $!A_n$ bukan premis minor \Elim{\lor} atau \Elim{\lexists};
  \item Untuk $i = 1$, \dots,~$n-1$, $!A_i$ merupakan premis minor
  \Elim{\lor} atau \Elim{\lexists}, dan $!A_{i+1}$ merupakan kesimpulan
  inferensi tersebut.
\end{enumerate}
\emph{Panjang} segmen tersebut ialah~$n$.
\end{defn}

Tidak setiap segmen merupakan sasaran eliminasi (jalan memutar). Segmen yang
menjadi sasaran disebut \emph{potongan}.

\begin{defn}
Suatu \emph{potongan} adalah segmen dengan $!A_n$ sebagai premis utama
!!a{elimination}, dan dengan $n=1$ serta $!A_1 = !A_n$ sebagai kesimpulan
!!a{introduction} atau~\FalseInt, ataupun dengan~$n>1$.
\end{defn}

Perhatikan bahwa segmen potongan tidak harus dimulai dengan kesimpulan
!!a{introduction}. Yang disyaratkan hanyalah bahwa segmen itu berakhir pada
premis utama !!a{elimination}. Akan tetapi, agar dihitung sebagai potongan,
segmen berpanjang~$1$ harus merupakan kesimpulan !!a{introduction} atau
\FalseInt.

\begin{defn}
\emph{Peringkat} suatu potongan ialah~$\depth{!A}$. \emph{Peringkat
potongan}~$\cutrank{\delta}$ dari !!a{proof}~$\delta$ adalah peringkat
maksimum potongan dalam~$\delta$. Suatu potongan \emph{maksimal} dalam
$\delta$ jika peringkatnya~$\cutrank{\delta}$, yakni jika peringkatnya
maksimum. \emph{Panjang potongan}~$\delta$ adalah jumlah panjang semua
potongan maksimalnya.
\end{defn}

\end{document}
